\begin{mistake}{2026-06-17}{05}

\begin{problemcontent}
设 \(f(x)\) 在 \(x=0\) 的某邻域内有定义，则 \(F(x) = f(x)|\sin x|\) 在 \(x=0\) 处可导的充分必要条件是（  ）

A. \(\lim_{x\to0} f(x)\) 存在

B. \(\lim_{x\to0} f(x) = f(0)\)

C. \(f(x)\) 在 \(x=0\) 处可导

D. \(\lim_{x\to0^-} f(x)\) 与 \(\lim_{x\to0^+} f(x)\) 均存在，且 \(\lim_{x\to0^-} f(x) = -\lim_{x\to0^+} f(x)\)
\end{problemcontent}

\begin{solutioncontent}
利用导数定义分别计算左右导数。
首先因为 \(f(x)\) 在 \(x=0\) 有定义，所以 \(F(0) = f(0)|\sin 0| = 0\)。

右导数（当 \(x \to 0^+\) 时，\(|\sin x| = \sin x\)）：
\[
F'_+(0) = \lim_{x \to 0^+} \frac{f(x)|\sin x| - 0}{x} = \lim_{x \to 0^+} \frac{f(x)\sin x}{x} = \lim_{x \to 0^+} f(x)
\]

左导数（当 \(x \to 0^-\) 时，\(|\sin x| = -\sin x\)）：
\[
F'_-(0) = \lim_{x \to 0^-} \frac{f(x)|\sin x| - 0}{x} = \lim_{x \to 0^-} \frac{-f(x)\sin x}{x} = -\lim_{x \to 0^-} f(x)
\]

\(F(x)\) 在 \(x=0\) 处可导的充要条件为左右导数存在且相等，即 \(\lim_{x \to 0^+} f(x) = -\lim_{x \to 0^-} f(x)\)。
故选 D。
\end{solutioncontent}

\mistakeknowledge{
\begin{itemize}
  \item \textbf{核心公式/定理}：分段点（或绝对值符号变化点）可导的充要条件是左右导数存在且相等：\(F'_+(x_0) = F'_-(x_0)\)。
  \item \textbf{易错点}：误认为需要 \(f(x)\) 在 \(x=0\) 连续。实际上 \(F(0)=0\) 恒成立，导数定义中只依赖 \(f(x)\) 在 \(x \to 0\) 时的极限，与 \(f(0)\) 的取值无关。
  \item \textbf{关联知识}：处理含绝对值的函数在零点的求导问题时，切忌直接套用求导法则，必须使用导数定义去绝对值后计算极限。
\end{itemize}}

\end{mistake}
